% META: {"id": "TNSI-LANG-EX-027", "chapitre": "TNSI-LANGAGES-ET-PROGRAMMATION", "type_objet": "exercice", "capacites": ["TNSI-LANGAGES-ET-PROGRAMMATION-C6"], "capacites_codes": ["C6"], "methodes": ["M1"], "parcours": 2, "competences": ["programmer"], "duree_min": 8, "mode_creation": "ex_nihilo", "sources_inspiration": [], "fichier_tex": "chapitres/TNSI-LANGAGES-ET-PROGRAMMATION/exercices/TNSI-LANG-EX-027.tex", "corrige_tex": "chapitres/TNSI-LANGAGES-ET-PROGRAMMATION/corriges/TNSI-LANG-CO-027.tex", "status": "approved"}
% BEGIN-VERIFY
% def mots_longs_imperatif(mots, longueur_min):
%     resultat = []
%     for mot in mots:
%         if len(mot) >= longueur_min:
%             resultat.append(mot.upper())
%     return resultat
% def mots_longs_fonctionnel(mots, longueur_min):
%     return [mot.upper() for mot in mots if len(mot) >= longueur_min]
% mots = ["chat", "ordinateur", "cle", "programmation"]
% assert mots_longs_imperatif(mots, 6) == mots_longs_fonctionnel(mots, 6)
% assert mots_longs_fonctionnel(mots, 6) == ["ORDINATEUR", "PROGRAMMATION"]
% END-VERIFY
\begin{exercice}{TNSI-LANG-EX-003}{2}{20}
La fonction suivante, écrite en style \textbf{impératif}, sélectionne et met en majuscules les mots d'une liste dont la longueur atteint un minimum donné :
\begin{python}
def mots_longs_imperatif(mots, longueur_min):
    resultat = []
    for mot in mots:
        if len(mot) >= longueur_min:
            resultat.append(mot.upper())
    return resultat
\end{python}
\begin{enumerate}
  \item Réécrire cette fonction en style \textbf{fonctionnel}, sous le nom \lstinline{mots_longs_fonctionnel}, en une seule expression (liste en compréhension).
  \item Vérifier que les deux versions renvoient le même résultat sur \lstinline{["chat", "ordinateur", "cle", "programmation"]} avec \lstinline{longueur_min=6}.
\end{enumerate}
\end{exercice}
